\section{The LAMP Protocol}
\label{sec:construction}

Section~\ref{sec:tech-overview} reduced the original matrix-multiplication
relation to the Freivalds relation and then to sampled encoded checks.  This
section formalizes the relation checked by \(\protocol\).  The important point is
that \(\protocol\) is not a single arithmetic-circuit relation: the arithmetic
checks are proved by an inner CP-SNARK, while Merkle membership and CP-Link are
verified outside the circuit.

\subsection{The Composite LAMP Relation}
\label{sec:relation-chain}

Let \(\mat{A},\mat{B},\mat{C}\in\mathbb{F}^{k\times k}\), and let
\(\mathcal{C}:\mathbb{F}^k\to\mathbb{F}^n\) be a linear code with
relative distance \(\delta\).  The target matrix-multiplication relation is
\[
  \mathcal{R}_{\mathsf{orig}}
  =
  \{(\mat{A},\mat{B},\mat{C}) :
    \mat{A}\mat{B}=\mat{C}\}.
\]
The verifier does not receive the original matrices as public input.  Instead, the
prover commits to row-wise encodings of the matrices and proves that the
committed objects satisfy the multiplication relation.

The relation checked by \(\protocol\) is the composite relation
\[
  \mathcal{R}_{\mathsf{LAMP}}
  =
  \mathcal{R}_{\mathsf{LAMP}}^{\mathsf{on}}
  \land
  \mathcal{R}_{\mathsf{LAMP}}^{\mathsf{off}} .
\]
Here \(\mathcal{R}_{\mathsf{LAMP}}^{\mathsf{on}}\) is the online relation proved by the CP-SNARK after the Freivalds challenge, query positions, and code-check points are sampled.  The relation \(\mathcal{R}_{\mathsf{LAMP}}^{\mathsf{off}}\) is the offline relation: it captures the values fixed by commitments before the online queries are known, and the verifier checks it through Merkle openings and CP-Link proofs.
Sections~\ref{sec:snark-relation} and~\ref{sec:openings-and-cplink} define
these two relations explicitly.
The circuit checks field equations on sampled encoded values, while the auxiliary verifier checks that those values are the ones bound by the pre-query commitments.

\noindent\textbf{Using the original matrices in an application circuit.}
The standard construction avoids placing all matrix entries in the CP-SNARK
witness: the external roots bind the encoded matrices, and the CP-SNARK commits
only to the \(\mathcal{O}(tk)\) sampled values used online.  If an application
circuit already uses \(\mat{A},\mat{B},\mat{C}\), it can instead commit to all
\(\mathcal{O}(k^2)\) entries and reuse those witness variables in the
\(\protocol\) checks.  This removes the separate Merkle/CP-Link bridge to the
application circuit while preserving
\(\mathcal{O}(tk+E_{\mathsf{code}})\) multiplication-check constraints; binding
the commitment to the intended application data remains application-specific.

% \subsection{Commitment Layout and Public Challenges}
% \label{sec:commitment-challenge}
% The prover first encodes the matrices row-wise:
% {\footnotesize
% \(\widehat{\mat{A}}=\enc(\mat{A}),
%   \widehat{\mat{B}}=\enc(\mat{B}),
%   \widehat{\mat{C}}=\enc(\mat{C})\).
% }
% For each codeword position \(j\in[n]\), it packs the three encoded columns into
% {\footnotesize
% \(\vect{m}_{ABC,j}
%   \coloneqq
%   \big(
%     \widehat{\mat{A}}[*,j]\,
%     \|\,\widehat{\mat{B}}[*,j]\,
%     \|\,\widehat{\mat{C}}[*,j]
%   \big)
%   \in\mathbb{F}^{3k}\).
% }
% Using a Pedersen commitment key \(\ck_{ABC}\), the prover samples
% \(o_{ABC,j}\sample\mathbb{F}\) and computes
% {\footnotesize
% \(\cm_{ABC,j}
%   \gets
%   \pedC(\ck_{ABC},\vect{m}_{ABC,j};o_{ABC,j})\).
% }
% The matrix root is the Merkle commitment to these Pedersen commitments:
% {\footnotesize
% \(\rt_{ABC}
%   \gets
%   \merkleC((\cm_{ABC,j})_{j\in[n]};\bot)\).
% }
% The root \(\rt_{ABC}\) is fixed before the Freivalds challenge is derived.
% After receiving \(\rt_{ABC}\), the verifier samples \(r\sample\mathbb{F}\). 
% The prover sets {\footnotesize \(\vect{r}=(1,r,r^2,\ldots,r^{k-1})\)}
% and computes
% {\footnotesize
% \(\vect{x}=\vect{r}\mat{A},
%   \vect{y}=\vect{x}\mat{B},
%   \vect{z}=\vect{r}\mat{C}\).
% }
% It then encodes the intermediate vectors:
% {\footnotesize
% \(\widehat{\vect{x}}=\enc(\vect{x}),
%   \widehat{\vect{y}}=\enc(\vect{y}),
%   \widehat{\vect{z}}=\enc(\vect{z})\).
% }
% For each \(j\in[n]\), the prover packs the encoded vector symbols into
% {\footnotesize
% \(\vect{m}_{XYZ,j}
%   \coloneqq
%   \big(
%     \widehat{\vect{x}}[j]\,
%     \|\,\widehat{\vect{y}}[j]\,
%     \|\,\widehat{\vect{z}}[j]
%   \big)
%   \in\mathbb{F}^{3}\).
% }
% Using a Pedersen commitment key \(\ck_{XYZ}\), the prover samples
% \(o_{XYZ,j}\sample\mathbb{F}\), computes
% {\footnotesize
% \(\cm_{XYZ,j}
%   \gets
%   \pedC(\ck_{XYZ},\vect{m}_{XYZ,j};o_{XYZ,j})\),
% }
% and sends the vector-symbol root
% {\footnotesize
% \(\rt_{XYZ}
%   \gets
%   \merkleC((\cm_{XYZ,j})_{j\in[n]};\bot)\)
% }
% to the verifier.
% Only after receiving \(\rt_{XYZ}\) does the verifier sample the query tuple
% {\footnotesize \(I=(j_1,\ldots,j_t)\sample [n]^t\).}
% It also samples public code-check points
% {\footnotesize
% \(\alpha_x,\alpha_y,\alpha_z\sample\mathbb{F}\setminus D\),
% }
% where \(D\) is the
% code evaluation domain.  These points are used by the circuit to check that
% {\footnotesize
% \(\widehat{\vect{x}},\widehat{\vect{y}},\widehat{\vect{z}}\)
% }
% are encodings of
% the claimed intermediate vectors.  For the Reed--Solomon code used in our implementation, this check is the barycentric out-of-domain consistency test recalled in Appendix~\ref{app:rs-barycentric}. 
\subsection{Commitment Layout and Public Challenges}
\label{sec:commitment-challenge}

The prover first encodes the matrices row-wise:
{\footnotesize
\(\widehat{\mat{A}}=\enc(\mat{A}),
  \widehat{\mat{B}}=\enc(\mat{B}),
  \widehat{\mat{C}}=\enc(\mat{C})\).
}
For each codeword position \(j\in[n]\), it packs the three \(k\)-dimensional
interleaved symbols, equivalently the three encoded columns, into
{\footnotesize
\(\vect{m}_{ABC,j}
  \coloneqq
  \big(
    \widehat{\mat{A}}[*,j]\,
    \|\,\widehat{\mat{B}}[*,j]\,
    \|\,\widehat{\mat{C}}[*,j]
  \big)
  \in\mathbb{F}^{3k}\).
}
Using a Pedersen commitment key \(\ck_{ABC}\), the prover samples
\(o_{ABC,j}\sample\mathbb{F}\) and computes
{\footnotesize
\(\cm_{ABC,j}
  \gets
  \pedC(\ck_{ABC},\vect{m}_{ABC,j};o_{ABC,j})\).
}
The matrix root is the Merkle commitment to these Pedersen commitments:
{\footnotesize
\(\rt_{ABC}
  \gets
  \merkleC((\cm_{ABC,j})_{j\in[n]};\bot)\).
}
Here and below, \(\bot\) indicates that no additional randomness is used in the
Merkle tree, since the leaves are Pedersen commitments that already provide
hiding. The root \(\rt_{ABC}\) is fixed before the Freivalds challenge is
derived.

After receiving \(\rt_{ABC}\), the verifier independently samples
\(r,s\sample\mathbb{F}\).  The prover sets
{\footnotesize
\(\vect{r}=(1,r,r^2,\ldots,r^{k-1})\) and
\(\vect{s}=(1,s,s^2,\ldots,s^{k-1})\)}
and computes
{\footnotesize
\(\vect{x}=\vect{r}\mat{A},
  \vect{y}=\vect{x}\mat{B},
  \vect{z}=\vect{r}\mat{C},
  \vect{b}=\vect{s}\mat{B}\).
}
It then encodes the intermediate vectors:
{\footnotesize
\(\widehat{\vect{x}}=\enc(\vect{x}),
  \widehat{\vect{y}}=\enc(\vect{y}),
  \widehat{\vect{z}}=\enc(\vect{z}),
  \widehat{\vect{b}}=\enc(\vect{b})\).
}
The encoded intermediate vectors form a \(4\)-fold interleaved codeword.  For
each \(j\in[n]\), the prover packs its \(j\)-th interleaved symbol into
{\footnotesize
\(\vect{m}_{XYZB,j}
  \coloneqq
  \big(
    \widehat{\vect{x}}[j]\,
    \|\,\widehat{\vect{y}}[j]\,
    \|\,\widehat{\vect{z}}[j]\,
    \|\,\widehat{\vect{b}}[j]
  \big)
  \in\mathbb{F}^{4}\).
}
Using a Pedersen commitment key \(\ck_{XYZB}\), the prover samples
\(o_{XYZB,j}\sample\mathbb{F}\), computes
{\footnotesize
\(\cm_{XYZB,j}
  \gets
  \pedC(\ck_{XYZB},\vect{m}_{XYZB,j};o_{XYZB,j})\),
}
and sends the vector-symbol root
{\footnotesize
\(\rt_{XYZB}
  \gets
  \merkleC((\cm_{XYZB,j})_{j\in[n]};\bot)\)
}
to the verifier.

The fourth vector supplies an independent fold of the committed encoding of
\(\mat{B}\); its soundness role is analyzed in Section~\ref{sec:security}.

Only after receiving \(\rt_{XYZB}\) does the verifier sample the query tuple
{\footnotesize \(I=(j_1,\ldots,j_t)\sample [n]^t\).}
It also samples public code-check points
{\footnotesize
\(\alpha_x,\alpha_y,\alpha_z,\alpha_b\sample\mathbb{F}\setminus D\),
}
where \(D\) is the
code evaluation domain. These points are used by the circuit to check that
{\footnotesize
\(\widehat{\vect{x}},\widehat{\vect{y}},\widehat{\vect{z}},
\widehat{\vect{b}}\)
}
are encodings of
the claimed intermediate vectors. For the Reed--Solomon code used in our
implementation, this check is the barycentric out-of-domain consistency test
recalled in Appendix~\ref{app:rs-barycentric}.


\subsection{The Online Relation}
\label{sec:snark-relation}
The online relation is the arithmetic relation proved inside the
inner CP-SNARK.  Its role is not to verify Merkle paths or Pedersen openings.
Instead, it checks only the field equations that certify the sampled encoded
positions are consistent with the Freivalds reduction.

The public statement contains the Merkle roots and verifier-sampled
challenges:
\[
\stmt = \big( \rt_{ABC},\rt_{XYZB},r,s,I,\alpha_x,\alpha_y,\alpha_z,\alpha_b
  \big).
\]
The roots are included in the statement to bind the CP-SNARK proof to the same
transcript and offline checks; the online field equations themselves do not
evaluate Merkle paths.

The complete SNARK witness is
  $\wit=(\comwit,\omega)$,
where the committed part of the witness is
\[
\comwit = \big( (\vect{m}_{ABC,j})_{j\in I},(\vect{m}_{XYZB,j})_{j\in I} \big),
\]
and the remaining private witness is
  $\omega
  =
  \big(
    \vect{x},\vect{y},\vect{z},\vect{b},
    \widehat{\vect{x}},\widehat{\vect{y}},\widehat{\vect{z}},
    \widehat{\vect{b}}
  \big).$
The CP-SNARK treats \(\comwit\) as the committed part of the witness.  The
prover first computes the SNARK-side witness commitments
\(\widetilde{\cm}_{ABC}\) and \(\widetilde{\cm}_{XYZB}\) using
\(\bPi_{\cp}.\com\), and then proves the online relation with witness
\(\wit=(\comwit,\omega)\).  These witness commitments are not checked against
the Merkle roots inside the circuit.  Instead, they are connected to the
externally committed Merkle leaves by the offline relation in
Section~\ref{sec:openings-and-cplink}.
For the CP-Link statements below, let \(\ck_{\cp,T}\) denote the public
commitment key used by the CP-SNARK witness commitment
\(\widetilde{\cm}_{T}\), for \(T\in\{ABC,XYZB\}\).
Write
  $\mathsf{EncCheck}_{\mathcal{C}}(\widehat{\vect{v}},\vect{v};\alpha)=1$
when the out-of-domain code-consistency check accepts that
\(\widehat{\vect{v}}\) is consistent with \(\enc(\vect{v})\) at the public
point \(\alpha\).  The online relation is defined in
Figure~\ref{relation:online}.

\begin{figure}[t]
\begin{tcolorbox}[colback=white, colframe=black]
\pseudocodeblock[mode=text]{
$\circ$ $\relation_{\mathsf{LAMP}}^{\mathsf{on}}(\stmt; \wit):$ \\
    $\qquad$parse \\
    $\qquad \qquad \stmt =
      \big(
        \rt_{ABC},\rt_{XYZB},r,s,I,\alpha_x,\alpha_y,\alpha_z,\alpha_b
      \big)$ \\
    $\qquad \qquad \wit=(\comwit,\omega)$ \\
    $\qquad \qquad \comwit =
      \big(
        (\vect{m}_{ABC,j})_{j\in I},
        (\vect{m}_{XYZB,j})_{j\in I}
      \big)$ \\
    $\qquad \qquad \omega =
      \big(
        \vect{x},\vect{y},\vect{z},\vect{b},
        \widehat{\vect{x}},\widehat{\vect{y}},\widehat{\vect{z}},
        \widehat{\vect{b}}
      \big)$ \\[0.5ex]
    $\qquad \vect{r}\gets(1,r,\ldots,r^{k-1})$ \\[0.5ex]
    $\qquad \vect{s}\gets(1,s,\ldots,s^{k-1})$ \\[0.5ex]
    $\qquad \mathsf{EncCheck}_{\mathcal{C}}
      (\widehat{\vect{x}},\vect{x};\alpha_x)=1$ \\
    $\qquad \mathsf{EncCheck}_{\mathcal{C}}
      (\widehat{\vect{y}},\vect{y};\alpha_y)=1$ \\
    $\qquad \mathsf{EncCheck}_{\mathcal{C}}
      (\widehat{\vect{z}},\vect{z};\alpha_z)=1$ \\
    $\qquad \mathsf{EncCheck}_{\mathcal{C}}
      (\widehat{\vect{b}},\vect{b};\alpha_b)=1$ \\[0.5ex]
    $\qquad \forall j\in I:$ \\
    $\qquad \qquad$parse
      $\vect{m}_{ABC,j}$ as
      $\big(
        \widehat{\mat{A}}[*,j]\,
        \|\,\widehat{\mat{B}}[*,j]\,
        \|\,\widehat{\mat{C}}[*,j]
      \big)$ \\
    $\qquad \qquad$parse
      $\vect{m}_{XYZB,j}$ as
      $\big(
        \widehat{\vect{x}}[j]\,
        \|\,\widehat{\vect{y}}[j]\,
        \|\,\widehat{\vect{z}}[j]\,
        \|\,\widehat{\vect{b}}[j]
      \big)$ \\[0.5ex]
    $\qquad \qquad
      \widehat{\vect{x}}[j]
      =
      \fold(\vect{r},\widehat{\mat{A}}[*,j])$ \\
    $\qquad \qquad
      \widehat{\vect{y}}[j]
      =
      \fold(\vect{x},\widehat{\mat{B}}[*,j])$ \\
    $\qquad \qquad
      \widehat{\vect{z}}[j]
      =
      \fold(\vect{r},\widehat{\mat{C}}[*,j])$ \\
    $\qquad \qquad
      \widehat{\vect{b}}[j]
      =
      \fold(\vect{s},\widehat{\mat{B}}[*,j])$ \\
    $\qquad \qquad
      \widehat{\vect{y}}[j]
      =
      \widehat{\vect{z}}[j]$
}
\end{tcolorbox}
\caption{The Online Relation}
\label{relation:online}
\end{figure}
\FloatBarrier


The encoding checks validate the four intermediate views, and the remaining
equations check the sampled folds and \(\widehat{\vect{y}}[j]=
\widehat{\vect{z}}[j]\).  Binding these sampled values to the pre-query roots
is the role of the offline relation.

\subsection{The Offline Relation}
\label{sec:openings-and-cplink}

For compactness, Figure~\ref{relation:lamp-off} writes the CP-Link statement
for \(T\in\{ABC,XYZB\}\) as
\[
  \stmt_{\mathsf{link}}^T
  \coloneqq
  \big(
    (\ck_{\cp,T},\widetilde{\cm}_{T}),
    (\ck_T,\cm_{T,j})_{j\in I}
  \big),
\]
where \(\ck_{\cp,T}\) is the witness-commitment key used by the CP-SNARK
commitment \(\widetilde{\cm}_{T}\), and \(\ck_T\) is the external Pedersen
commitment key for the Merkle leaves of type \(T\).

The offline relation connects the sampled witnesses used by the
CP-SNARK to the commitments that were fixed before the query positions were
sampled.  The term ``offline'' indicates that these checks are performed outside the SNARK circuit. In our system, it also reflects that the committed data are fixed before the online sampled checks are carried out.

For each queried position \(j\in I\), the verifier receives Merkle openings for the external Pedersen leaves $\cm_{ABC,j}$ and $\cm_{XYZB,j}$.
The Merkle openings certify that these leaves are contained in the roots
\(\rt_{ABC}\) and \(\rt_{XYZB}\), respectively.  In addition, the verifier checks
CP-Link proofs connecting the CP-SNARK witness commitments
\(\widetilde{\cm}_{ABC}\) and \(\widetilde{\cm}_{XYZB}\) to the corresponding
external Pedersen leaves.

For each \(T\in\{ABC,XYZB\}\), CP-Link proves that the SNARK-side commitment
contains the ordered concatenation
\(\vect{m}_{T,j_1}\|\cdots\|\vect{m}_{T,j_t}\) committed by the corresponding
external leaves.  Appendix~\ref{app:cplink-instant} gives the concrete
blockwise relation.


The offline relation is defined in Figure~\ref{relation:lamp-off}.

\begin{figure}[!t]
\begin{tcolorbox}[colback=white, colframe=black]
\pseudocodeblock[mode=text]{
$\circ\;
\relation_{\mathsf{LAMP}}^{\mathsf{off}}\big($
    $
      \rt_{ABC},\rt_{XYZB},I,
      \widetilde{\cm}_{ABC},\widetilde{\cm}_{XYZB},
      $ \\
    $\qquad
      \pi_{\mathsf{link}}^{ABC},\pi_{\mathsf{link}}^{XYZB},$ \\
    $\qquad
      (\cm_{ABC,j},\pi_{\mathsf{mem},j}^{ABC})_{j\in I},
      (\cm_{XYZB,j},\pi_{\mathsf{mem},j}^{XYZB})_{j\in I}
    \big):$ \\[0.5ex]
    $\qquad \forall j\in I:$ \\
    $\qquad \qquad
        \merkleV(
            \rt_{ABC},
            j,
            \cm_{ABC,j},
            \pi_{\mathsf{mem},j}^{ABC}
        ) = 1$ \\
    $\qquad \qquad
        \merkleV(
            \rt_{XYZB},
            j,
            \cm_{XYZB,j},
            \pi_{\mathsf{mem},j}^{XYZB}
        ) = 1$ \\[0.5ex]
    $\qquad
        \bPi_{\mathsf{link}}.\verify(
            \para_{\mathsf{link}},
            \stmt_{\mathsf{link}}^{ABC},
            \pi_{\mathsf{link}}^{ABC}
        ) = 1$ \\
    $\qquad
        \bPi_{\mathsf{link}}.\verify(
            \para_{\mathsf{link}},
            \stmt_{\mathsf{link}}^{XYZB},
            \pi_{\mathsf{link}}^{XYZB}
        ) = 1$
}
\end{tcolorbox}
\caption{The Offline Relation}
\label{relation:lamp-off}
\end{figure}
\FloatBarrier

Together, CP-Link and Merkle membership establish
\(\widetilde{\cm}_{T}\leftrightarrow(\cm_{T,j})_{j\in I}
\leftrightarrow\rt_T\) for \(T\in\{ABC,XYZB\}\).  Hence the CP-SNARK checks
values bound by the pre-query roots without verifying Pedersen openings or
Merkle paths inside the circuit.

% \begin{table}[t]
% \centering
% \footnotesize
% \setlength{\tabcolsep}{4pt}
% \begin{tabularx}{\columnwidth}{@{}lcc@{}}
% \toprule
% \textbf{Component} & \textbf{Prover} & \textbf{Verifier} \\
% \midrule
% ABC encoding
% & \(E_{\enc}(k,n)\)
% & -- \\

% ABC commit
% & \(\mathcal{O}(3nk)\,\mathbb{G}+\mathcal{O}(n)\,\mathbb{H}\)
% & -- \\

% Intermediate
% & \(\mathcal{O}(k^2)\,\mathbb{F}\)
% & -- \\

% XYZ encoding
% & \(E_{\enc}(k,n)\)
% & -- \\

% XYZ commit
% & \(\mathcal{O}(n)\,\mathbb{G}+\mathcal{O}(n)\,\mathbb{H}\)
% & -- 
% \\
% Merkle openings
% & \(\mathcal{O}(t\log n) \mathbb{H}\)
% & \(\mathcal{O}(t\log n)\mathbb{H} \) 
% \\
% % CP-SNARK
% % & \(
% % \mathcal{O}(tk)+E_{\mathsf{code}} +
% % E_{\mathsf{snark}}^\mathcal{P}
% % \)
% % & \(E_{\mathsf{snark}}^\mathcal{V}\) 
% CP-SNARK
% & \(\begin{array}{c}
% \mathcal{O}(tk)+E_{\mathsf{code}}\ \text{constraints}\\
% E_{\mathsf{snark}}^\mathcal{P}
% \end{array}\)
% & \(E_{\mathsf{snark}}^\mathcal{V}\)
% \\
% CP-Link
% & \(E_{\mathsf{link}}^\mathcal{P}(k,t)\)
% & \(E_{\mathsf{link}}^\mathcal{V}(k,t)\) \\
% \bottomrule
% \end{tabularx}
% \vspace{6pt}
% \caption{Asymptotic costs for \(\protocol\).}
% \label{tab:cost-summary}
% \end{table}

\subsection{Protocol Algorithms and Complexity}
\label{sec:protocol-complexity}
Figure~\ref{fig:protocol} presents the public-coin version of
\(\protocol\).  The protocol can be made non-interactive via the
Fiat--Shamir transform, by deriving \((r,s)\) after \(\rt_{ABC}\) and deriving
\(I,\alpha_x,\alpha_y,\alpha_z,\alpha_b\) after \(\rt_{XYZB}\).

The proof consists of the CP-SNARK proof, the sampled leaf commitments and
Merkle openings, and the CP-Link proofs.  The sampled matrix-column blocks and
vector-symbol blocks are private CP-SNARK witnesses; they are not serialized as
public proof material.
The public statement contains the roots
\(\rt_{ABC},\rt_{XYZB}\) and the verifier-sampled challenges. \\

For each \(T\in\mathcal{T}=\{ABC,XYZB\}\), the CP-Link statement contains
the SNARK-side witness commitment and the \(t\) corresponding external leaf
commitments; the witness contains their messages and opening randomness.\\

% \paragraph{Cost model}
\noindent\textbf{Cost model.}
We report the cost of \protocol.  The computation of the claimed output matrix \(C=AB\) is outside this cost model.
We do count the work performed by the \protocol prover after the matrices are available, including encoding, commitment generation, computation of the Freivalds intermediate witnesses, sampled Merkle openings, CP-SNARK proving, and CP-Link proving.

Let \(E_{\enc}(k,n)\) denote the cost of the corresponding encoding phase.
Let \(E_{\mathsf{code}}=E_{\mathsf{code}}(k,n)\) denote the constraint cost of the intermediate
code-consistency checks; for the Reed--Solomon check in
Appendix~\ref{app:rs-barycentric}, this is \(\mathcal{O}(k+n)\) after
domain-dependent coefficients are precomputed.
We write \(\mathbb{G}\) for group operations used by Pedersen commitments,
\(\mathbb{F}\) for field operations, and \(\mathbb{H}\) for hash evaluations.
Let \(E_{\mathsf{snark}}^{\mathcal{P}}\) and \(E_{\mathsf{snark}}^{\mathcal{V}}\) denote the CP-SNARK proving and
verification costs, respectively. Let \(E_{\mathsf{link}}^\mathcal{P}(k,t)\) and \(E_{\mathsf{link}}^\mathcal{V}(k,t)\) denote the prover and verifier costs, respectively, of the selected CP-Link construction.
For the QALink instantiation in Appendix~\ref{app:cplink-instant}, each of the
two link statements has \(q=t\) external commitments, so
\(E_{\mathsf{link}}^\mathcal{V}(k,t)\) includes \(2(t+2)\) Miller-loop inputs,
batched into one final exponentiation, plus \(\mathcal{O}(t)\) group scalar
multiplications for batching.

Table~\ref{tab:cost-summary} summarizes the asymptotic costs at the level of the main protocol components. \\

\begin{table}[ht]
\centering
\footnotesize
\setlength{\tabcolsep}{4pt}
\begin{tabularx}{\columnwidth}{@{}lcc@{}}
\toprule
\textbf{Component} & \textbf{Prover} & \textbf{Verifier} \\
\midrule
ABC encoding
& \(E_{\enc}(k,n)\)
& -- \\

ABC commit
& \(\mathcal{O}(3nk)\,\mathbb{G}+\mathcal{O}(n)\,\mathbb{H}\)
& -- \\

Intermediate
& \(\mathcal{O}(k^2)\,\mathbb{F}\)
& -- \\

XYZB encoding
& \(E_{\enc}(k,n)\)
& -- \\

XYZB commit
& \(\mathcal{O}(n)\,\mathbb{G}+\mathcal{O}(n)\,\mathbb{H}\)
& -- 
\\
Merkle openings
& \(\mathcal{O}(t\log n) \mathbb{H}\)
& \(\mathcal{O}(t\log n)\mathbb{H} \) 
\\
% CP-SNARK
% & \(
% \mathcal{O}(tk)+E_{\mathsf{code}} +
% E_{\mathsf{snark}}^\mathcal{P}
% \)
% & \(E_{\mathsf{snark}}^\mathcal{V}\) 
CP-SNARK
& \(\begin{array}{c}
\mathcal{O}(tk)+E_{\mathsf{code}}\ \text{constraints}\\
E_{\mathsf{snark}}^\mathcal{P}
\end{array}\)
& \(E_{\mathsf{snark}}^\mathcal{V}\)
\\
CP-Link
& \(E_{\mathsf{link}}^\mathcal{P}(k,t)\)
& \(E_{\mathsf{link}}^\mathcal{V}(k,t)\) \\
\bottomrule
\end{tabularx}
\vspace{6pt}
\caption{Asymptotic costs for \(\protocol\).}
\label{tab:cost-summary}
\end{table}


\noindent\textbf{Prover work.}
The prover first encodes the three input matrices and commits to the packed
encoded matrix columns.  The ABC commitment phase consists of \(n\) Pedersen
commitments to blocks of length \(3k\), followed by a Merkle root over
the resulting leaves.  After the independent challenges \((r,s)\) are sampled, the
prover computes the intermediate witnesses
  $\vect{x}=\vect{r}\mat{A},
  \vect{y}=\vect{x}\mat{B},
  \vect{z}=\vect{r}\mat{C},
  \vect{b}=\vect{s}\mat{B},$
which costs \(\mathcal{O}(k^2)\) field operations.  This is distinct from the
excluded native computation of \(C=AB\).  The prover then encodes
\(\vect{x},\vect{y},\vect{z},\vect{b}\), commits to their encoded symbols, opens the
sampled Merkle paths, proves the online CP-SNARK relation, and produces the
CP-Link proofs.

For each queried position \(j\in I\), the CP-SNARK circuit checks four
length-\(k\) fold equations:
\[
\begin{aligned}
  &\fold(\vect{r},\widehat{\mat{A}}[*,j]),\quad
  \fold(\vect{x},\widehat{\mat{B}}[*,j]),\\
  &\fold(\vect{r},\widehat{\mat{C}}[*,j]),\quad
  \fold(\vect{s},\widehat{\mat{B}}[*,j]).
\end{aligned}
\]
Thus the sampled arithmetic checks contribute \(\mathcal{O}(tk)\) constraints. 

The intermediate-vector code-consistency checks add
\(E_{\mathsf{code}}\) constraints. For the fixed-rate
Reed--Solomon setting used in the experiments, \(n=2k\), so this extra cost is
linear in \(k\) and the total online relation remains linear in \(k\) for fixed
\(t\). \\

% \paragraph{Verifier work}
\noindent\textbf{Verifier work.}
The verifier does not encode matrices, compute intermediate vectors, or build
commitment roots.  Its work consists of CP-SNARK verification, verification of
the sampled Merkle authentication paths, and CP-Link verification.  Therefore,
% using the notation above, the verifier cost is $E_{\mathsf{snark}}^\mathcal{V} + \mathcal{O}(t\log n)\mathbb{H} + E_{\mathsf{link}}^\mathcal{V}(k,t).$ \\
using the notation above, the verifier cost is
$E_{\mathsf{snark}}^\mathcal{V} + \mathcal{O}(t\log n)\mathbb{H} +
E_{\mathsf{link}}^\mathcal{V}(k,t)$, where the QALink verifier contributes
\(\mathcal{O}(t)\) pairing terms. \\

% \paragraph{Comparison}
\noindent\textbf{Comparison.}
A naive SNARK arithmetization of matrix multiplication costs
\(\mathcal{O}(k^3)\) constraints, while a Freivalds-based SNARK baseline computes
the vector-matrix products inside the circuit and costs \(\mathcal{O}(k^2)\) constraints.  
% \(\protocol\) moves the vector-matrix products to native prover work and checks only sampled encoded-column relations inside the CP-SNARK.  
% As a result, the main circuit constraint count is \(\mathcal{O}(tk)\).
\(\protocol\) moves the vector-matrix products to native prover work and checks
sampled encoded-column relations plus intermediate code-consistency equations
inside the CP-SNARK.
As a result, the main circuit constraint count is
\(\mathcal{O}(tk+E_{\mathsf{code}})\), which is linear in \(k\) for the
fixed-rate code used in our implementation.


\begin{figure*}[p]
\centering
\scriptsize
\begingroup
\setlength{\fboxsep}{2pt}
\resizebox{0.98\textwidth}{!}{%
\fbox{%
\procedure[]{\protocol}{%
    \mathcal{P}(\mat{A},\mat{B},\mat{C}) \< \< \mathcal{V} \\
    \widehat{\mat{A}},\widehat{\mat{B}},\widehat{\mat{C}} \gets \enc(\mat{A}),\enc(\mat{B}),\enc(\mat{C}) \in \mathbb{F}^{k \times n} \< \< \\
    \text{for each } j\in[n]: \\ 
    \quad \vect{m}_{ABC,j} \coloneqq \big( \widehat{\mat{A}}[*,j]\, \|\,\widehat{\mat{B}}[*,j]\, \|\,\widehat{\mat{C}}[*,j] \big) \in\mathbb{F}^{3k} \\
    \quad o_{ABC,j}\sample\mathbb{F} \\
    \quad \cm_{ABC,j} \gets \pedC(\ck_{ABC},\vect{m}_{ABC,j};o_{ABC,j}) 
    \\
    \rt_{ABC} \gets \merkleC((\cm_{ABC,j})_{j\in[n]};\bot) 
    \< \sendmessageright*[1.3cm]{\rt_{ABC}} 
    \\
    \< \sendmessageleft*[1.3cm]{r,s} \<  r,s \sample \FF 
    \\
    \vect{r}\coloneqq(1,r,\ldots,r^{k-1}) \\ 
    \vect{s}\coloneqq(1,s,\ldots,s^{k-1}) \\
    \vect{x}\gets\vect{r}\mat{A}, \quad \vect{y}\gets\vect{x}\mat{B}, \quad \vect{z}\gets\vect{r}\mat{C}, \quad \vect{b}\gets\vect{s}\mat{B}
    \\ \widehat{\vect{x}},\widehat{\vect{y}},\widehat{\vect{z}},\widehat{\vect{b}} \gets \enc(\vect{x}),\enc(\vect{y}),\enc(\vect{z}),\enc(\vect{b}) \\ 
    \text{for each } j\in[n]: \\ 
    \quad \vect{m}_{XYZB,j} \coloneqq \big( \widehat{\vect{x}}[j]\, \|\,\widehat{\vect{y}}[j]\, \|\,\widehat{\vect{z}}[j]\, \|\,\widehat{\vect{b}}[j] \big) \in\mathbb{F}^{4} \\ 
    \quad o_{XYZB,j}\sample\mathbb{F} \\ 
    \quad \cm_{XYZB,j} \gets \pedC(\ck_{XYZB},\vect{m}_{XYZB,j};o_{XYZB,j}) \\ 
    \rt_{XYZB} \gets \merkleC((\cm_{XYZB,j})_{j\in[n]};\bot) \< \sendmessageright*[1.3cm]{\rt_{XYZB}} \< \\
    \< \< I=(j_1,\ldots,j_t)\sample[n]^t 
    \\ 
    \< \sendmessageleft*[1.3cm]{I, \alpha_x, \alpha_y, \alpha_z, \alpha_b} \< \alpha_x,\alpha_y,\alpha_z,\alpha_b\sample\mathbb{F}\setminus D 
    \\
    \mathcal{T}\coloneqq(ABC,XYZB) \\
    \text{For every } T\in\mathcal{T}\text{ and every } j\in I,\\ 
    \quad \pi_{\mathsf{mem},j}^{T} \gets \merkleO((\cm_{T,\ell})_{\ell\in[n]},j) \\[0.3ex]
    \stmt \coloneqq
      \big(\rt_{ABC},\rt_{XYZB},r,s,I,\alpha_x,\alpha_y,\alpha_z,\alpha_b\big) \\
    \comwit \coloneqq
      \big((\vect{m}_{T,j})_{j\in I}\big)_{T\in\mathcal{T}}, \quad
    \omega \coloneqq
      \big(
      \vect{x},\vect{y},\vect{z},\vect{b},
      \widehat{\vect{x}},\widehat{\vect{y}},\widehat{\vect{z}},\widehat{\vect{b}}
      \big) \\
    \wit \coloneqq \big(\comwit,\omega\big) \\
    \text{For every } T\in\mathcal{T},\\
    \quad \widetilde{\cm}_{T} \gets
      \bPi_{\cp}.\com(\para,(\vect{m}_{T,j})_{j\in I};\widetilde{o}_{T}) \\
    \stmt_{\cp} \gets
      \big(
      \stmt,
      (\widetilde{\cm}_{T})_{T\in\mathcal{T}}
      \big) \\
    \pi_{\cp} \gets \bPi_{\cp}.\prove(\para,\stmt_{\cp},\wit) \\
    \text{For every } T\in\mathcal{T},\\
    \quad \pi_{\mathsf{link}}^{T} \gets
      \bPi_{\mathsf{link}}.\prove(
      \para_{\mathsf{link}},
      \stmt_{\mathsf{link}}^T,
      \wit_{\mathsf{link}}^T) \\
    \pi_{\protocol}
        \coloneqq
        \big(
        \pi_{\cp},
        (\widetilde{\cm}_{T})_{T\in\mathcal{T}},
        \{(\cm_{T,j},\pi_{\mathsf{mem},j}^{T})\}_{T\in\mathcal{T},j\in I},
(\pi_{\mathsf{link}}^{T})_{T\in\mathcal{T}}
        \big)  \< \sendmessageright*[1.3cm]{\pi_\protocol}  \<
   \text{Parse } \pi_{\protocol} 
   \\
    \< \< \stmt_{\cp} \gets \big( \stmt,(\widetilde{\cm}_{T})_{T\in\mathcal{T}} \big) \\
    \< \< b_{\cp} \gets \bPi_{\cp}.\verify \big( \para, \stmt_{\cp}, \pi_{\cp}) \\
    \< \< \text{For every } T\in\mathcal{T}: \\
    \< \< \qquad \text{For every } j\in I: \\
    \< \< \qquad \qquad b_{\mathsf{mem},j}^{T} \gets
      \merkleV( \rt_{T}, j, \cm_{T,j}, \pi_{\mathsf{mem},j}^{T}) \\
    \< \< \qquad b_{\mathsf{link}}^{T} \gets
      \bPi_{\mathsf{link}}.\verify \big(
      \para_{\mathsf{link}},
      \stmt_{\mathsf{link}}^T,
      \pi_{\mathsf{link}}^{T}
      \big) \\
    \< \< \text{Accept iff }
      b_{\cp}
      \land
      \bigwedge_{T\in\mathcal{T}}
      \left(
        b_{\mathsf{link}}^{T}
        \land
        \bigwedge_{j\in I} b_{\mathsf{mem},j}^{T}
      \right)
}
}
}
\endgroup
\caption{The \protocol\ protocol}
\label{fig:protocol}
\end{figure*}
